% !TeX root = booklet-01-preview.tex
\section{نمونه آموزشی: ولتاژ و جریان}
\label{sec:electrical-basics}
هدف این نمونه آن است که بتوانیم عملکرد یک \lr{CMOS Inverter} را از روی
اتصال‌های مدار توضیح دهیم. ابتدا کمیت‌های لازم را می‌شناسیم و سپس از یک مدل ساده
برای ترانزیستور استفاده می‌کنیم. بخش مقدماتی این نمونه، توضیح تکمیلی برای ورود
به مدل \lr{Switch} در بخش \secnum{۱.۳} کتاب است؛ مثال‌های عددی آن، داده فناوری یا نتیجه شبیه‌سازی نیستند.

\subsection{بار و جریان چه تفاوتی دارند؟}
بار الکتریکی را با \(Q\) و واحد \lr{coulomb} با نماد \(\mathrm{C}\) نمایش می‌دهیم.
جریان، آهنگ عبور بار از یک مقطع است. اگر در مدت \(\Delta t\)، بار خالص
\(\Delta Q\) در جهت مرجع از مقطع عبور کند، جریان متوسط برابر است با:
\begin{equation}\label{eq:average-current}
I_{\mathrm{avg}}=\frac{\Delta Q}{\Delta t}.
\end{equation}
علامت \(\Delta\) یعنی «تغییر» یا «اختلاف»؛ در اینجا \(\Delta t\) طول بازه زمانی است.
کسر در رابطه \eqref{eq:average-current} یعنی مقدار بار را بر مدت عبور آن تقسیم می‌کنیم.
واحد جریان \lr{ampere} است و \(\qty{1}{\ampere}=\qty{1}{\coulomb\per\second}\).
جهت جریان قراردادی، جهت حرکت بار مثبت است؛ حرکت \lr{Electron} در فلز در جهت مخالف آن است.

\textbf{مثال آموزشی.} فرض کنید در مدت \(\qty{2}{\nano\second}\)، بار خالص
\(\qty{6}{\pico\coulomb}\) در جهت مرجع از یک مقطع عبور کند.
پیشوندهای \lr{pico} و \lr{nano} به‌ترتیب به معنی \(10^{-12}\) و \(10^{-9}\) هستند.
جای‌گذاری، همراه با تبدیل واحد، می‌دهد:
\begin{align}
I_{\mathrm{avg}}
&=\frac{6\times10^{-12}\,\mathrm{C}}{2\times10^{-9}\,\mathrm{s}}
=\frac{6}{2}\times10^{-12-(-9)}\,\mathrm{A}\notag\\
&=3\times10^{-3}\,\mathrm{A}=\qty{3}{\milli\ampere}.
\end{align}
در تقسیم توان‌های هم‌پایه، توان مخرج از توان صورت کم می‌شود؛ بنابراین
\(-12-(-9)=-3\). این نتیجه جریان متوسط است؛ از آن نمی‌توان ثابت‌بودن جریان
در تمام لحظه‌های بازه را نتیجه گرفت.

\subsection{ولتاژ همیشه نسبت به یک مرجع است}
ولتاژ، اختلاف پتانسیل الکتریکی است؛ از دید انرژی، اختلاف پتانسیل دو نقطه نشان می‌دهد
به‌ازای هر واحد بار مثبت، انرژی پتانسیل چقدر تفاوت دارد.
واحد آن \lr{volt} است: \(\qty{1}{\volt}=\qty{1}{\joule\per\coulomb}\).
اگر پتانسیل دو نقطه \(a\) و \(b\) نسبت به یک مرجع، \(V_a\) و \(V_b\) باشد، تعریف می‌کنیم:
\begin{equation}\label{eq:voltage-difference}
V_{ab}=V_a-V_b.
\end{equation}
پس اگر \(V_a=\qty{1}{\volt}\) و \(V_b=\qty{0.4}{\volt}\)،
داریم \(V_{ab}=\qty{0.6}{\volt}\)، ولی \(V_{ba}=\qty{-0.6}{\volt}\).
عوض‌کردن ترتیب پایانه‌ها علامت را تغییر می‌دهد؛ این نکته در تعریف \(V_{gs}\) مهم خواهد بود.

\clearpage
\subsection{مرجع، اتصال و شرط عبور جریان}
\lr{Ground} گره‌ای است که پتانسیل آن را در تحلیل صفر در نظر می‌گیریم.
هر \lr{Node}، مجموعه نقاط متصل با سیم ایدئال است که پتانسیل یکسان دارند.
\lr{Ground} لزوماً اتصال فیزیکی به زمین نیست؛ صفر بودن ولتاژ یک گره نیز
به معنی نبود بار یا نبود جریان در تمام شاخه‌های متصل به آن نیست.

وجود ولتاژ به‌تنهایی برای عبور جریان از هر مسیر دلخواه کافی نیست؛ مسیر باید امکان
هدایت داشته باشد. در مدل ایدئال، \lr{Open Circuit} مسیر هدایت را قطع می‌کند،
درحالی‌که \lr{Short Circuit} اختلاف ولتاژ دو سر خود را صفر می‌کند.
جریان اتصال کوتاه را سایر اجزای مدار تعیین می‌کنند؛ نمی‌توان گفت همیشه صفر یا همیشه بی‌نهایت است.

برای یک مقاومت مثبت \(R\)، جریان مرجع \(I\) را از سر با علامت مثبت ولتاژ \(V\)
به سر منفی انتخاب می‌کنیم. در مدل مقاومت خطی، \lr{Ohm's Law} می‌گوید:
\begin{equation}\label{eq:ohm}
V=RI\quad\Longrightarrow\quad I=\frac{V}{R}.
\end{equation}
برای رسیدن به عبارت دوم، دو طرف عبارت اول را بر \(R\) تقسیم کرده‌ایم.
واحد مقاومت \lr{ohm} است. اگر \(V=\qty{1}{\volt}\) و
\(R=\qty{1}{\kilo\ohm}=\qty{1000}{\ohm}\)، جریان
\(I=\qty{0.001}{\ampere}=\qty{1}{\milli\ampere}\) خواهد بود.
این قانون مدل مقاومت است؛ هنوز مجاز نیستیم آن را مدل دقیق ترانزیستور بدانیم.

\section{از ولتاژ به صفر و یک منطقی}
مقدار منطقی با مقدار ولتاژ یکسان نیست. برای نمونه، در مدار آموزشی با تغذیه
\(V_{\mathrm{DD}}=\qty{1}{\volt}\)، ولتاژ نزدیک صفر را نماینده صفر منطقی و
ولتاژ نزدیک تغذیه را نماینده یک منطقی می‌گیریم.
عدد یک منطقی واحد ندارد؛ \(\qty{1}{\volt}\) یک کمیت فیزیکی است.
تغذیه انتخاب‌شده فقط برای مثال است و مشخصات فناوری خاصی را بیان نمی‌کند.

در مدار واقعی، صفر و یک با بازه‌های مجاز ولتاژ تعریف می‌شوند؛ هر ولتاژ میانی
الزاماً یک مقدار منطقی معتبر نیست. در این نمونه فقط دو ورودی دقیقاً برابر
صفر و \(V_{\mathrm{DD}}\) را تحلیل می‌کنیم. تحلیل بازه‌ها و تحمل نویز در جزوه سوم می‌آید.
این شیوه نمایش سطوح با معرفی کتاب سازگار است \cite[بخش \secnum{۱.۳}، صفحه چاپی ۸]{weste}.

\textbf{نکته کلیدی.} صفر بودن ورودی یعنی ولتاژ آن نسبت به \lr{Ground} صفر است؛
این عبارت به معنی قطع‌بودن همه ترانزیستورهای مدار نیست. نوع ترانزیستور و اختلاف
ولتاژ پایانه‌های آن تعیین می‌کنند کدام مسیر هدایت دارد.

\clearpage
\section{ترانزیستور به‌عنوان \lr{Switch}}
\subsection{پایانه کنترل را از مسیر هدایت جدا کنیم}
ترانزیستور \lr{MOSFET} چهار پایانه \lr{Gate}، \lr{Source}، \lr{Drain} و
\lr{Body} دارد. در مدل مقدماتی، ولتاژ \lr{Gate} وضعیت مسیر میان
\lr{Source} و \lr{Drain} را کنترل می‌کند. جریان مسیر هدایت از خود
\lr{Gate} عبور نمی‌کند؛ در مدل ایدئال، عایق \lr{Gate} جریان رسانشی آن را صفر می‌کند.
بااین‌حال، برای تغییر ولتاژ \lr{Gate} باید ظرفیت آن شارژ یا دشارژ شود؛
پس این مدل، نبود هر نوع جریان ورودی در هنگام تغییر را ثابت نمی‌کند
\cite[بخش \secnum{۱.۳}]{weste}\cite[اسلایدهای ۵ تا ۸]{mos}.

ولتاژ کنترل \lr{nMOS} را \(V_{gs}=V_g-V_s\) می‌نامیم؛ بنابراین
فقط دانستن \(V_g\) کافی نیست. در \lr{Inverter} این نمونه،
\lr{Source} ترانزیستور \lr{nMOS} به صفر و \lr{Source} ترانزیستور \lr{pMOS}
به \(V_{\mathrm{DD}}\) متصل است. بدنه‌ها نیز به‌ترتیب به همین ریل‌ها متصل‌اند.
فرض می‌کنیم تغذیه برای روشن‌شدن هر ترانزیستور در ورودی مناسب کافی است.

\subsection{قاعده کنترل در همین اتصال مشخص}
با اتصال‌های یادشده، ورودی بالا \lr{nMOS} را روشن و \lr{pMOS} را خاموش می‌کند؛
ورودی پایین نتیجه معکوس دارد. این همان مدل \lr{Switch} کتاب و اسلاید است
\cite[صفحه چاپی ۸]{weste}\cite[اسلاید ۲۱]{intro}.
در این مدل، \lr{ON} یعنی مسیر بسته و \lr{OFF} یعنی مسیر باز.
\textbf{روشن بودن، وجود مسیر را بیان می‌کند؛ جریان غیرصفر را تضمین نمی‌کند.}
اگر اختلاف ولتاژ لازم وجود نداشته باشد، مسیر روشن نیز می‌تواند جریان صفر داشته باشد.

جدول \ref{tab:switch} فقط برای دو ورودی ریل‌به‌ریل و اتصال‌های همین نمونه است؛
قاعده عمومی برای همه مدارهای \lr{Pass Transistor} نیست.
\begin{table}[ht]
\centering
\caption{وضعیت مدل \lr{Switch} در اتصال مخصوص \lr{Inverter}}
\label{tab:switch}
\begin{tabular}{ccc}
\toprule
ورودی & \lr{nMOS} & \lr{pMOS}\\
\midrule
صفر منطقی & قطع & وصل\\
یک منطقی & وصل & قطع\\
\bottomrule
\end{tabular}
\end{table}

\textbf{حد اعتبار.} مقاومت مسیر روشن، \lr{Leakage} و روند زمانی تغییر وضعیت
در این مدل حذف شده‌اند. روابط جریان و ناحیه‌های عملکرد در ادامه جزوه اول،
و اثرهای \lr{Nonideal} در جزوه دوم آموزش داده خواهند شد.

\clearpage
\section{ساخت و تحلیل \lr{CMOS Inverter}}
\subsection{ابتدا اتصال‌ها را بخوانیم}
\lr{Complementary Metal-Oxide-Semiconductor (CMOS)} از دو نوع مکمل
ترانزیستور استفاده می‌کند. در \lr{Inverter} شکل \ref{fig:inverter}،
\lr{Gate} هر دو ترانزیستور ورودی مشترک \(A\) دارد و اتصال مشترک
\lr{Drain}ها خروجی \(Y\) است. ترانزیستور بالایی مسیر \lr{Pull-up}
از تغذیه به خروجی و ترانزیستور پایینی مسیر \lr{Pull-down} از خروجی به زمین را می‌سازد
\cite[بخش \secnum{۱.۴.۱}، صفحه چاپی ۹]{weste}.

\begin{figure}[ht]
\centering
\begin{latin}
\begin{circuitikz}[american,scale=1.05]
\draw (0,3) node[pmos] (p) {};
\draw (0,0) node[nmos] (n) {};
\draw (p.S) -- (0,4.2) node[above] {$V_{\mathrm{DD}}$};
\draw (n.S) -- (0,-1.2) node[ground] {};
\draw (p.D) -- (0,1.5) -- (n.D);
\draw (0,1.5) to[short,*-o] (2,1.5) node[right] {$Y$};
\draw (p.G) -- (-1.7,3) -- (-1.7,0) -- (n.G);
\draw (-1.7,1.5) to[short,*-o] (-2.6,1.5) node[left] {$A$};
\node[right] at (0.6,3) {pMOS};
\node[right] at (0.6,0) {nMOS};
\end{circuitikz}
\end{latin}
\caption{بازترسیم آموزشی \lr{CMOS Inverter} بر پایه شکل \secnum{۱.۱۱} کتاب و اسلاید ۲۳
\cite{weste,intro}. بدنه‌ها در نماد ساده نمایش داده نشده‌اند؛ اتصال آن‌ها در متن مشخص است.}
\label{fig:inverter}
\end{figure}

سیم ورودی به \lr{Gate}ها می‌رود و از مسیر خروجی جداست.
در سمت چپ شکل، نقطه اتصال نشان می‌دهد هر دو \lr{Gate} ولتاژ یکسان می‌گیرند.
حباب کوچک کنار \lr{Gate} ترانزیستور بالا نشانه \lr{pMOS} است.
خط‌های اتصال، سیم ایدئال‌اند؛ از روی شکل نباید مسیر رسانشی مستقیمی از \(A\) به \(Y\) فرض کرد.

\clearpage
\subsection{ورودی صفر: چرا خروجی یک می‌شود؟}
فرض می‌کنیم \(V_A=\qty{0}{\volt}\).
\lr{nMOS} خاموش است، زیرا \(V_{gs,n}=V_A-0=0\).
\lr{pMOS} روشن است، زیرا اختلاف ولتاژ کنترل مناسب آن
\(V_{sg,p}=V_{\mathrm{DD}}-V_A=V_{\mathrm{DD}}\) است.
پس مسیر خروجی به زمین قطع و مسیر آن به تغذیه وصل می‌شود.
در حالت ماندگار مدل ایدئال و بدون بار رسانشی خروجی، نتیجه می‌گیریم:
\begin{equation}\label{eq:inverter-low}
V_A=0\quad\Longrightarrow\quad V_Y=V_{\mathrm{DD}}.
\end{equation}
اگر خروجی ابتدا ولتاژ کمتری داشته باشد، ظرفیت آن از مسیر بالایی شارژ می‌شود؛
نتیجه \eqref{eq:inverter-low} مقدار نهایی است و زمان رسیدن به آن را تعیین نمی‌کند.

\subsection{ورودی یک: چرا خروجی صفر می‌شود؟}
اکنون \(V_A=V_{\mathrm{DD}}\). برای \lr{pMOS} داریم
\(V_{sg,p}=V_{\mathrm{DD}}-V_A=0\)، پس خاموش است.
برای \lr{nMOS} داریم \(V_{gs,n}=V_{\mathrm{DD}}\)، پس روشن است.
خروجی از تغذیه جدا و به زمین متصل می‌شود. در همان مدل و حالت ماندگار:
\begin{equation}\label{eq:inverter-high}
V_A=V_{\mathrm{DD}}\quad\Longrightarrow\quad V_Y=0.
\end{equation}
اگر ظرفیت خروجی قبلاً شارژ شده باشد، از مسیر پایینی دشارژ می‌شود.
جدول \ref{tab:inverter} دو نتیجه را به زبان منطق بیان می‌کند.

\begin{table}[ht]
\centering
\caption{جدول درستی \lr{Inverter} در مدل ایدئال و حالت ماندگار}
\label{tab:inverter}
\begin{tabular}{cc}
\toprule
ورودی \(A\) & خروجی \(Y\)\\
\midrule
۰ & ۱\\
۱ & ۰\\
\bottomrule
\end{tabular}
\end{table}

عملکرد منطقی با \(Y=\overline{A}\) نمایش داده می‌شود؛ خط بالای \(A\)
علامت \lr{NOT} است و به معنی منفی عددی نیست. این جدول را از اتصال‌ها استخراج کردیم؛
صرفاً نام گیت را حفظ نکردیم.

\clearpage
\section{کنترل فهم و مرز نتیجه}
در هر دو حالت پایدار بررسی‌شده، یکی از دو ترانزیستور قطع است؛ بنابراین
در مدل ایدئال مسیر رسانشی پیوسته‌ای از تغذیه به زمین نداریم.
این نتیجه به معنی صفر بودن توان واقعی مدار در همه شرایط نیست:
هنگام تغییر ورودی، ظرفیت‌ها شارژ و دشارژ می‌شوند، ممکن است دو ترانزیستور
برای مدتی هم‌زمان هدایت کنند و در مدار واقعی \lr{Leakage} نیز وجود دارد.
این محدودیت‌ها مانع استفاده آموزشی از مدل ساده نمی‌شوند؛ محدوده نتیجه را مشخص می‌کنند.

\subsection{پرسش‌ها و پاسخ‌های تشریحی}
\begin{enumerate}
\item \textbf{اگر ورودی صفر باشد، چرا خروجی هم صفر نیست؟}
ورودی به \lr{Gate}ها وصل است، نه مستقیماً به خروجی.
صفر شدن ورودی، ترانزیستور بالا را روشن می‌کند و خروجی به تغذیه می‌رسد.
\item \textbf{آیا ترانزیستور روشن همیشه جریان می‌دهد؟}
خیر. روشن بودن امکان هدایت را فراهم می‌کند. پس از رسیدن خروجی بی‌بار
به ریل مربوط، جریان ماندگار مدل ایدئال می‌تواند صفر شود.
\item \textbf{چرا در تعریف ولتاژ کنترل، ولتاژ \lr{Source} را کم می‌کنیم؟}
کنترل با اختلاف پتانسیل انجام می‌شود، نه با ولتاژ مطلق \lr{Gate}.
برای مثال، اگر \(V_g=\qty{1}{\volt}\) و \(V_s=\qty{0.4}{\volt}\)،
آنگاه \(V_{gs}=\qty{0.6}{\volt}\)، نه \(\qty{1}{\volt}\).
روشن بودن به مقایسه این اختلاف با آستانه نیاز دارد؛ مقدار آستانه در این مثال داده نشده است.
\item \textbf{اگر هر دو مسیر قطع شوند، آیا خروجی الزاماً صفر است؟}
خیر. خروجی می‌تواند \lr{Floating} باشد و بار قبلی را روی ظرفیت خود نگه دارد.
مدل ساده بدون شرایط اولیه، ولتاژ یکتایی تعیین نمی‌کند.
اگر هر دو مسیر روشن شوند نیز یک مسیر از تغذیه به زمین ایجاد می‌شود؛
برای محاسبه جریان و ولتاژ آن باید مدل دقیق‌تری داشت
\cite[صفحه چاپی ۱۰]{weste}.
\item \textbf{آیا این نمونه زمان تغییر خروجی یا رفتار ورودی میانی را محاسبه کرد؟}
خیر. فقط حالت ماندگار برای دو ورودی ریل‌به‌ریل تحلیل شد.
برای زمان به ظرفیت و جریان، و برای ورودی میانی به روابط جریان و ناحیه عملکرد نیاز داریم.
\end{enumerate}

پس از این نمونه، مسیر مطالعه نسخه کامل از قوانین مدار و خازن به گیت‌های پایه و
ساختار \lr{MOS} ادامه می‌یابد. سپس روابط جریان استخراج می‌شوند تا دلیل محدودیت‌های
مدل \lr{Switch} نیز قابل تحلیل باشد.
